\section{Теоретични основи на библиотеката и моделите за съхранение}
\label{sec:theoretical_background}

Теоретичната основа на настоящата разработка не се свежда до една езикова възможност или до един тип база данни. Библиотеката комбинира идеи от теорията за обектно-релационното съпоставяне (ОРС), статичното типизиране в C++, изграждането на междинни представяния по време на компилация и сравнителния анализ на различни модели за съхранение. За да бъде разбираема по-късната архитектура, е необходимо първо да се отдели логическият слой на описване на данните от физическия слой на тяхното специфично материализиране за дадения свързващ компонент.

В контекста на тази дипломна работа терминът обектно-релационно съпоставяне по време на компилация не трябва да се разбира стеснено само като SQL техника. По-точно той обозначава архитектурен подход, при който моделът на данните, формата на заявките и част от ограниченията върху операциите се описват чрез типовата система на езика, а не чрез конфигурация по време на изпълнение. Това позволява една и съща логическа структура да бъде пренесена към релационни, документни, ключ-стойност, графови и широки колонни системи, без специфичните за свързващия компонент различия да се скриват или подменят.

\subsection{Таксономия на проблемната област}

Преди разглеждане на отделните компоненти е полезно да се фиксира общата картина на проблемната област. От една страна стои C++ моделът, в който приложението описва типове същности, полета, връзки и ограничения. От друга страна стоят различни семейства системи за съхранение, всяко със собствен модел на физическа организация, достъп и оптимизация. Между тях е разположен типизираният слой за междинно представяне на заявките, който служи като посредник между логиката на приложението и конкретния слой на конектора.

\begin{figure}[H]
\centering
\begin{tikzpicture}[node distance=1.5cm and 1.4cm]
    \node[ormaccent] (cpp) {\textbf{C++ логически модел}\\типове същности, полета, връзки};
    \node[ormblock, below=of cpp] (ir) {\textbf{Типизиран междинно представяне на заявките}\\полета, правила, заместители};
    \node[ormblock, below left=of ir] (rel) {Релационен\\модел};
    \node[ormblock, below=of ir] (doc) {Документен\\модел};
    \node[ormblock, below right=of ir] (kv) {Ключ-стойност\\модел};
    \node[ormblock, below=of doc, xshift=-3.2cm] (graph) {Графов\\модел};
    \node[ormblock, below=of doc, xshift=3.2cm] (wide) {Широк колонен\\модел};

    \draw[ormline] (cpp) -- (ir);
    \draw[ormline] (ir) -- (rel);
    \draw[ormline] (ir) -- (doc);
    \draw[ormline] (ir) -- (kv);
    \draw[ormline] (ir) -- (graph);
    \draw[ormline] (ir) -- (wide);
\end{tikzpicture}
\caption{Таксономия на логическия модел, междинното представяне на заявките и семействата свързващи компоненти}
\label{fig:storage_taxonomy}
\end{figure}

Фигура~\ref{fig:storage_taxonomy} подчертава централното теоретично допускане на библиотеката: общото между свързващите компоненти не е техният синтаксис, а логическият модел на данните и операциите. Следователно валидната абстракция трябва да бъде разположена по-високо от SQL, BSON, Redis команди, Cypher или CQL. Тя трябва да работи на равнището на полета, предикати, проекции, връзки и допустими операции.

\subsection{Обектно-релационното съпоставяне като архитектурен подход}

Класическото ОРС свързва обектния модел на приложението с неговата персистентна схема~\cite{fowler2002}. Във варианта по време на компилация това свързване се конструира не чрез отражение (отражение) и конфигурации по време на изпълнение, а чрез типове, шаблони и \code{constexpr} стойности. Практически това означава, че изборът на полета, допустимите операции и част от съвместимостта със свързващия компонент се проверяват от компилатора. Грешки като обръщение към несъществуващо поле или използване на съединяване (\code{JOIN}) върху конектор без \code{supports\_joins} престават да бъдат изненади по време на изпълнение.

Съществената особеност е, че ОРС по време на компилация не означава липса на данни по време на изпълнение. Стойностите на параметрите се подават в момента на изпълнение, но структурата на заявката е вече фиксирана и валидирана. Така се получава ясно разграничение между \emph{логическа форма} и \emph{стойности по време на изпълнение}. Това разграничение е ключово за тезата, защото позволява едновременно статични гаранции и реална работа с външни бази данни.

От теоретична гледна точка този подход доближава ОРС слоя до компилаторен конвейер. Приложният код описва семантична структура, компонентите по време на компилация я превръщат в междинно представяне, а едва най-долният слой материализира тази структура в специфичен език или протокол на свързващия компонент. Затова по-късните архитектурни решения в библиотеката могат да се анализират през понятия като междинно представяне, статична проверка и материализация в целевата система.

\subsection{Статичен модел на данните в C++}

В разглежданата библиотека домейн моделът се описва чрез C++ агрегати, чиито полета са от тип \code{property<T, Name>}. По този начин всяко поле носи едновременно стойностен тип и символно име на колоната или атрибута. Допълнително \code{table\_name\_trait<T>} позволява типът да бъде свързан с логически контейнер за съхранение --- таблица, колекция, префикс на ключ, label или друго специфично за компонента понятие.

Тази организация има важно теоретично следствие: C++ типът става първичен носител на логическата схема. Отделните свързващи компоненти не налагат нов потребителски модел, а само интерпретират един и същ модел по различен начин. Именно затова по-късно може да се говори за еднакъв логически обект, който в релационен компонент се превръща в таблица, а в документен --- в документ с вложени или референтни полета.

\begin{table}[H]
\centering
\small
\caption{Съпоставка между C++ конструкциите и теоретичните ОРС понятия}
\label{tab:cpp_to_orm_concepts}
\begin{tabularx}{\textwidth}{L{3.3cm}L{4.0cm}Y}
\toprule
\textbf{Конструкция} & \textbf{Теоретична роля} & \textbf{Практически ефект в библиотеката} \\
\midrule
\code{property<T, Name>} & атрибут от логическата схема & носи тип на стойността и символно име за материализация \\
\code{relationship<...>} & логическа връзка между обекти & отделя връзката от физическия механизъм за съхранение \\
\code{table\_name\_trait<T>} & логически контейнер & служи като име на таблица, колекция, префикс на ключ или label \\
\code{field<\&T::m>} & адресиране на атрибут в пространството на заявката & свързва слоя на същностите с междинното представяне на заявките \\
\code{connector\_trait<DB>} & материализация в целевата система & определя как логическият модел се изпълнява физически \\
\bottomrule
\end{tabularx}
\end{table}

Следователно ОРС слоят не започва от драйвер или от рендиране на SQL. Той започва от статично описание на смисъла на данните. Това е и причината слоят на същностите да бъде фундаментален за цялата архитектура: ако този слой е неясен или непълен, няма как по-долу да се получи коректна и надграждаема материализация.

\begin{lstlisting}[language=C++,caption={Извадка от \code{property} и \code{relationship} в слоя на същностите},label={lst:property_relationship}]
template <typename T, detail::string_literal Name>
struct property
{
    using value_type = T;

    [[nodiscard]] static constexpr std::string_view column_name() noexcept
    {
        return Name.view();
    }

    T value{};
};

enum class store_as { reference, embed };

template <store_as Strategy, typename T, detail::string_literal Name>
struct relationship
{
    static constexpr store_as strategy = Strategy;
    using related_type = T;

    [[nodiscard]] static constexpr std::string_view field_name() noexcept
    {
        return Name.view();
    }
};
\end{lstlisting}

Listing~\ref{lst:property_relationship} показва, че слоят на същностите е минималистичен по форма, но достатъчно богат по смисъл. \code{property} не е просто контейнер за стойност; той носи и име, а \code{relationship} не е просто член-данна; тя фиксира стратегията за материализация като част от типа.

\begin{figure}[H]
\centering
\begin{tikzpicture}[node distance=1.4cm and 1.3cm]
    \node[ormaccent] (entity) {\textbf{Тип същност}\\агрегат от \code{property} и \code{relationship}};
    \node[ormblock, below left=of entity] (props) {Скаларни\\полета};
    \node[ormblock, below right=of entity] (rels) {Логически\\връзки};
    \node[ormblock, below=of props] (schema) {Логическа\\схема};
    \node[ormblock, below=of rels] (physical) {Физическа\\структура};
    \node[ormnote, below=1.2cm of $(schema)!0.5!(physical)$, minimum width=7.8cm] (note)
        {Един и същ C++ модел може да се материализира като SQL таблица,
         документ, префикс на ключ + стойност, label + свойства или широк колонен ред.};

    \draw[ormline] (entity) -- (props);
    \draw[ormline] (entity) -- (rels);
    \draw[ormline] (props) -- (schema);
    \draw[ormline] (rels) -- (schema);
    \draw[ormline] (schema) -- (physical);
\end{tikzpicture}
\caption{Преминаване от C++ описание на същност към логическа схема и физическа материализация}
\label{fig:entity_to_schema_mapping}
\end{figure}

\subsection{Типизирано междинно представяне на заявките}

Вместо заявката да се описва директно като SQL или друго текстово представяне, библиотеката използва типизирано междинно представяне на заявкитете. Неговите елементи са \code{field}, \code{Rule}, \code{placeholder} и обекти на заявките като \code{select\_query}, \code{insert\_query}, \code{update\_query} и \code{delete\_query}. Този модел дава две предимства. Първо, структурата на заявката може да бъде анализирана по време на компилация и да участва в \code{static\_assert} проверки. Второ, едно и също междинно представяне може да бъде рендирано към различни езици и протоколи на свързващите компоненти.

Теоретично междинното представяне играе ролята на междинен език между домейн модела и конкретния слой за изпълнение. Това е аналогично на междинните представяния в компилаторите: потребителят работи с високо ниво на абстракция, а свързващият компонент получава специализирана форма, оптимизирана за собствения си модел. Именно поради това междинното представяне е логически по-близо до семантиката на операциите, отколкото до синтаксиса на конкретна заявка.

\begin{lstlisting}[language=C++,caption={Извадка от \code{select\_query} и композицията на \code{where} клаузи},label={lst:select_query_excerpt}]
template <typename Response, typename Joins, typename Wheres,
          typename Limits, typename Groups, typename Orders>
struct select_заявка : select_заявка_tag
{
    using response_type = Response;
    using row_type = projected_type<Response>;

    template <typename RuleType>
        requires is_rule<RuleType>
    [[nodiscard]] constexpr auto where(RuleType r) const
    {
        using NewWheres = detail::append_type_t<Wheres, RuleType>;
        return select_заявка<Response, Joins, NewWheres, Limits, Groups, Orders>(
            response_, joins_,
            detail::кортеж_cat(wheres_, detail::orm_кортеж<RuleType>(r)),
            limits_, groups_, orders_);
    }
};
\end{lstlisting}

Listing~\ref{lst:select_query_excerpt} илюстрира важно теоретично свойство: конструкторът на заявки не мутира текст, а изгражда нови типизирани обекти. Това означава, че последователността от операции \code{select(...).where(...).group\_by(...)} може да бъде разглеждана като конструиране на нов тип, а не като поредица от странични ефекти върху низ.

\begin{figure}[H]
\centering
\begin{tikzpicture}[node distance=1.35cm and 1.1cm]
    \node[ormaccent] (заявка) {\textbf{Потребителска форма на заявка}\\\code{select}/\code{insert}/\code{update}/\code{deleteq}};
    \node[ormblock, below=of заявка] (field) {\code{field}, \code{Rule}, \code{placeholder}};
    \node[ormblock, below=of field] (typed) {Типизирано\\междинно представяне\\структура и ограничения};
    \node[ormblock, below left=of typed] (render) {Рендиране\\към целеви език};
    \node[ormblock, below right=of typed] (execute) {Изпълнение и\\свързване (свързване) на параметри};

    \draw[ormline] (заявка) -- (field);
    \draw[ormline] (field) -- (typed);
    \draw[ormline] (typed) -- (render);
    \draw[ormline] (typed) -- (execute);
\end{tikzpicture}
\caption{Конвейер на типизираното междинно представяне от приложния интерфейс до изпълнението в целевата система}
\label{fig:query_ir_pipeline}
\end{figure}

\subsection{Теория на връзките и стратегии за съхранение}

Връзките между обекти са един от най-трудните аспекти при обобщаване на обектно-релационното съпоставяне извън релационния свят. В релационен контекст една връзка често се материализира чрез външен ключ и съединяване (\code{JOIN}). В документен контекст обаче същата логическа връзка може да се реализира като вграден документ, а в ключ-стойност контекст --- чрез вторично търсене или денормализиране.

Библиотеката абстрахира тези варианти чрез \code{relationship<Strategy, T, Name>} и \code{store\_as::reference}/\code{store\_as::embed}. \code{reference} обозначава логическа външна препратка, която свързващият компонент може да реализира чрез външен ключ, търсене, отделен ключ или графова връзка. \code{embed} обозначава физическо вграждане на свързаните данни в същата физическа структура. По този начин връзката се описва веднъж логически, а интерпретацията остава задача на конектора.

\begin{figure}[H]
\centering
\begin{tikzpicture}[node distance=1.5cm and 1.4cm]
    \node[ormaccent] (rel) {\textbf{Логическа връзка}\\същност A $\leftrightarrow$ същност B};
    \node[ormblock, below left=of rel] (reference) {\code{store\_as::reference}\\външен ключ / търсене / ребро};
    \node[ormblock, below right=of rel] (embed) {\code{store\_as::embed}\\вложен документ / денормализирани данни};
    \node[ormnote, below=3cm of rel, minimum width=8.1cm] (conclusion)
        {Изборът е част от логическото описание, но цената и конкретната реализация
         зависят от свързващия компонент.};

    \draw[ormline] (rel) -- (reference);
    \draw[ormline] (rel) -- (embed);
\end{tikzpicture}
\caption{Двете основни стратегии за материализация на връзки}
\label{fig:relationship_strategies}
\end{figure}

\begin{table}[H]
\centering
\small
\caption{Сравнение между стратегиите \code{reference} и \code{embed}}
\label{tab:relationship_strategies}
\begin{tabularx}{\textwidth}{L{2.8cm}L{3.2cm}L{3.0cm}Y}
\toprule
\textbf{Стратегия} & \textbf{Логическа идея} & \textbf{Естествени целеви системи} & \textbf{Основен компромис} \\
\midrule
\code{reference} & отделна идентичност на свързания обект & SQL, графови системи, търсене по ключ & необходим е допълнителен механизъм за навигация или свързване \\
\code{embed} & физическо включване в родителската структура & документни системи, денормализирани записи & по-трудно независимо обновяване и споделяне на данни \\
\bottomrule
\end{tabularx}
\end{table}

\subsection{Релационен модел}

Релационният модел организира данните в таблици, редове и колони. Основните му предимства са ясна схема, декларативен език за заявки и добре дефинирани ограничения върху целостта. За системите за обектно-релационно съпоставяне това е естествената среда: свойствата на C++ обекта лесно се съпоставят с колони, а връзките --- с външни ключове. Операции като съединяване (\code{JOIN}), групиране (\code{GROUP BY}) и транзакции са част от нормалния езиков модел.

Именно релационният свързващ компонент е естественият базов сценарий, спрямо който може да се оцени до каква степен една библиотека за съпоставяне запазва удобството си и извън SQL света. В настоящата разработка този базов сценарий по-късно се реализира основно чрез \code{SQLiteDB}, а преносимостта на междинното представяне се анализира допълнително чрез \code{PostgreSQLDB} и \code{MySQLDB}.

\subsection{Документен модел}

Документният модел организира данните в колекции от документи, чиито полета могат да бъдат вложени и не е задължително всички документи да имат еднаква физическа форма. Това го прави естествено съвместим със стратегията за вграждане (\code{embed}) и по-гъвкав при денормализирани структури. От друга страна, твърде директното пренасяне на релационни абстракции върху документен свързващ компонент може да бъде подвеждащо. Не всяка релационна операция има еквивалент със същата цена и семантика.

В контекста на настоящата работа документният модел е важен като тест за това дали междинното представяне описва логика, а не синтаксис на конкретен език. Ако едно и също междинно представяне може да се интерпретира като BSON-подобен филтър и проекция, без потребителят да пише платформено-специфични документни заявки, тогава действително е постигнат слой на абстракция, независим от съхранението.

\subsection{Ключ-стойност модел}

Ключ-стойност системите организират достъпа около ключ и стойност. Те са особено ефективни, когато моделът на достъп е ясен и концентриран около търсене по идентификатор. Недостатъкът е, че общите заявки, базирани на предикати, и релационните операции са силно ограничени. Затова една абстракция за съпоставяне върху ключ-стойност компонент не може да обещава същата изразителност като върху SQL система.

Този модел е полезен за разграничаване между \emph{логически обект} и \emph{допустим модел на достъп}. Обектът може да бъде един и същ, но не всяка заявка е смислена върху всяка физическа организация. В хранилището това по-късно се вижда ясно при \code{RedisDB}, където логиката за търсене по ключ е естествена, а стилът, базиран на съединявания, е архитектурно неуместен.

\subsection{Графов модел}

Графовият модел представя данните като възли, ребра и свойства. Силата му е в изразяването на връзки като първокласни обекти, а не като производни от външни ключове. Това го прави подходящ за сценарии със сложни навигации и операции за обхождане. В същото време графовият модел налага собствен език на заявки и собствена интуиция за селекция и филтриране.

За библиотеката това е показателен сценарий, защото демонстрира, че \code{connector\_trait} може да пренесе общо междинно представяне към \code{MATCH}/\code{RETURN}/\code{WHERE} структура, без потребителят да пише Cypher директно. Ако това е възможно, значи общият слой действително е разположен над конкретния синтактичен слой.

\subsection{Широк колонен модел}

Широките колонни системи, каквато е Cassandra, се основават на разпределени таблици с ключове за разделяне (разделяне) и групиране (групиране). Макар синтактично да напомнят релационни таблици, те имат различна теория на достъпа: правилната заявка е тази, която следва физическия модел на ключовото разпределение. Следователно наличието на поле в C++ модела не означава автоматично, че може да бъде използвано свободно във \code{WHERE} клаузи.

Точно този тип ограничения оправдава подхода с възможности и ограничения на библиотеката. Не е достатъчно междинното представяне да е типово коректно; то трябва да е и допустимо спрямо модела на съхранение. По тази причина широкият колонен модел е особено ценен за теоретичната защита на ограниченията по време на компилация.

\begin{table}[H]
\centering
\small
\caption{Съпоставка между основните модели за съхранение}
\label{tab:storage_models}
\begin{tabularx}{\textwidth}{L{2.3cm}L{2.7cm}L{2.8cm}Y}
\toprule
\textbf{Модел} & \textbf{Единица за съхранение} & \textbf{Типични връзки} & \textbf{Основно ограничение} \\
\midrule
Релационен & таблица/ред & външен ключ, съединяване & необходимост от съгласувана схема и силна типова дисциплина \\
Документен & документ & вграждане или препратка & не всяка релационна операция е естествена или евтина \\
Ключ-стойност & ключ и стойност & търсене по ключ, хеш полета & силно ограничени заявки с предикати извън ключовия достъп \\
Графов & възел и ребро & връзки за обхождане & специфичен език за заявки и графова семантика \\
Широк колонен & разделен (разделянеed) ред & ключово-ориентирани зависимости & заявките трябва да следват логиката на разделянето \\
\bottomrule
\end{tabularx}
\end{table}

\subsection{Тестова опора на теоретичните твърдения}

Теоретичните твърдения в тази глава не са откъснати от кода. Слоят на същностите се валидира пряко в \path{tests/модулни/test_property.cpp} и \path{tests/модулни/test_relationship.cpp}, където се проверяват имената на колоните, разпознаването на типовете същности, различаването между \code{store\_as::reference} и \code{store\_as::embed}, както и логиката за извеждане (извеждане) на отношения едно-към-много. Типизираното междинно представяне е покрито в \path{tests/модулни/test_заявка.cpp}, където се демонстрират предикати, базирани на полета, форми на съединяване, композиция на обновявания и семантика на изтриване. Допълнително \path{tests/модулни/test_cpp26_отражение.cpp} показва, че описанието на полетата по време на компилация може да бъде надградено и чрез подход с отражение, когато инструментариумът го позволява.

Тази връзка между теория и тестова опора е съществена. Тя показва, че представените понятия не са свободни академични метафори, а реални архитектурни единици, които могат да бъдат проверявани в изходния код и в автоматизираните тестове.

\subsection{Граници на унифицираната абстракция}

След изложеното може да се направи важен извод. Унифицираната абстракция не означава, че всички свързващи компоненти стават еднакви. Унифицирането е възможно на равнище логически модел, поле, връзка, междинно представяне и базови операции. Отвъд тази граница започват различията: SQL \code{JOIN} няма универсален смисъл в Redis; логиката за вграждане няма идентична цена в класическа SQL таблица; обхождането на граф не се свежда до обикновен \code{SELECT}; ограниченията, продиктувани от разделянето в Cassandra, не могат да бъдат маскирани като обикновени предикатни правила.

Поради това архитектурно правилният подход не е да се скрият всички различия, а да се направят явни и формално контролирани. Именно тази роля изпълняват \code{connector\_trait}, labelите за възможности и ограниченията по време на компилация, разгледани в следващите глави. Точно там ще стане видимо как теоретичната рамка от настоящата глава се превръща в работеща архитектура на библиотеката.
